% STATUS: written 2026-11-06 (English gate clean; sources read in full as recorded below).
% Sources read in full: Hasse--Schmidt 1937 (J. reine angew. Math. 177, 215--237; the four differentiation rules (1)--(4) on p. 218).
% Not yet read: Hasse 1936 (Crelle 175); stated here only as Hasse and Schmidt describe it.
\section[Higher derivations in characteristic $p$: Hasse and Schmidt 1937][Higher derivations]{Higher derivations in characteristic $p$: Hasse and Schmidt 1937}\label{sec:origins:higher-derivations}

In characteristic $p$ the usual higher derivatives also lose information. For
example, the $p$-th derivative of every polynomial is $0$. The theory of
higher derivations began as an answer to this problem. It began in the study
of algebraic function fields.

The solution came from a letter. The paper of Hasse and Schmidt~\cite{HasseSchmidt1937}
grew out of it, and it is written in German. We follow it in our translation. On 25 May 1936
F.~K.~Schmidt wrote to Hasse. He was teaching the recent theory of higher
differential quotients of Hasse. For his lectures, he had worked it out in a
different form. He wanted to introduce the quotients globally (``im
Gro\ss en'') from the start. He did not want to start from the power series
at a point~\cite[p.~215]{HasseSchmidt1937}.

Schmidt had assumed that his definition agreed with one given by
Teichm\"uller. Then the paper of Teichm\"uller appeared. It showed Schmidt
that the two definitions differ. Hasse published the result as a free and
detailed version of the letter. In some points he changed the argument of
Schmidt.

We now turn to the construction itself. The setting is a separably generated algebraic function field $K$ of one
variable. Its constant field $k$ is arbitrary. Let $x$ be a separating
element of $K/k$, and let $y \in K$. Let $\xi$ be a copy of $x$ in an
auxiliary field. The authors expand $y$ as a power series in powers of
$\xi - x$. The coefficients lie in $K$~\cite[\S 2, p.~217]{HasseSchmidt1937}.

\begin{definition}[{\cite[p.~217]{HasseSchmidt1937}}]\label{def:origins:hasse-schmidt}
The \emph{differential quotient of order $\nu$} of $y$ with respect to $x$,
written $\dfrac{d^\nu y}{dx^\nu}$, is the coefficient $y_\nu$ of
$(\xi - x)^\nu$ in this expansion ($\nu = 0, 1, 2, \ldots$).
\end{definition}

The authors stress what these quotients are. They are \translated{not the
exact equivalent of the differential quotients of analysis, but rather of
the Taylor coefficients}{Die so definierten Differentialquotienten sind nicht
das genaue \"Aquivalent zu den Differentialquotienten der Analysis, sondern
vielmehr zu den Taylorkoeffizienten}{German}
$\frac{1}{\nu!}\frac{d^\nu}{dx^\nu}$~\cite[p.~217]{HasseSchmidt1937}.

This is the main point of the construction. The division by $\nu!$ is part
of the definition. So the definition makes sense in characteristic~$p$.

The quotients depend only on $x$ and $y$. They do not depend on the
auxiliary choices. They satisfy four rules: a sum rule, a rule for
constants, a product rule and a power
rule~\cite[p.~218, (1)--(4)]{HasseSchmidt1937}. In the first order they give
the derivatives already used in the arithmetic of~$K/k$.

In modern language, the maps $D_\nu(y) = d^\nu y/dx^\nu$ form a \emph{higher derivation} of $K$
over $k$~\cite[p.~218]{HasseSchmidt1937}. It is also called a \emph{Hasse--Schmidt derivation}~\cite[p.~218]{HasseSchmidt1937}. It is a
sequence of $k$-linear maps $D_0 = \mathrm{id}, D_1, D_2, \ldots$ with
\[
D_\nu(yz) = \sum_{i+j=\nu} D_i(y)\, D_j(z) \qquad (\nu \ge 0).
\]
When $k$ contains the rational numbers, every derivation $D$ gives such a
sequence: $D_\nu = D^\nu/\nu!$. In characteristic $p$ we cannot divide by
$\nu!$. For this reason the maps $D_\nu$ must be given separately. The
algebraic theory of higher derivations is the subject of
Chapter~\ref{ch:higher-multi}.

The notion earned its place on three counts in the years that followed. Baer had
linked derivations to purely inseparable extensions of exponent one in 1927, and
Weisfeld generalized his results to extensions of any exponent
\cite[pp.~307--318]{Weisfeld1965Purely}. For algebras, Roy and Sridharan
introduce an equivalence among higher derivations and show that its classes
correspond to isomorphism classes of certain filtered modules
\cite[Introduction, pp.~115--116]{RoySridharan1968Higher}. And the notion leaves
the algebraic setting: Miller asks on a Banach algebra whether every higher
derivation has the form $\exp D$, and proves restricted converses
\cite[pp.~301--331]{Miller1970Higher}.

Heerema carried the notion into the theory of local rings. His survey of the
subject is written to show where higher derivations lead, and what they lead to
is the automorphism theory of complete local rings
\cite[Introduction, p.~1212]{Heerema1970Higher}.

So far, derivations have served the theory of fields. In the same years, Ritt
and his school made a field with a derivation an object of study in its own
right.
